% Abstrak ID & EN dan kata kunci

\section*{Abstrak}
\noindent
Matriks korelasi antarsubtes Ujian Tertulis Berbasis Komputer--Seleksi Nasional Berdasarkan Tes (UTBK-SNBT) 2026 tidak tersedia bagi publik. Statistik deskriptif per subtes yang beredar bersumber dari tabel internal Panitia Seleksi Nasional Penerimaan Mahasiswa Baru (SNPMB) yang tidak resmi dan belum terverifikasi, sehingga disebut parameter nasional tidak resmi. Penelitian ini bertujuan mengestimasi struktur korelasi kedelapan subtes dengan data \textit{crowdsourced} dari platform X ($N = 1.886$), sekaligus mengidentifikasi karakteristik bias pada metode itu serta implikasinya bagi estimasi peringkat nasional. Rerata korelasi \textit{off-diagonal} pada \textit{raw sample} sebesar $0,300$ tergolong moderat-rendah untuk serangkaian tes kemampuan kognitif yang saling berkaitan. Perbandingan dengan parameter nasional tidak resmi menunjukkan bias seleksi ke arah skor tinggi yang mengindikasikan \textit{restriction of range}. Koreksi Thorndike Case III berbasis iterasi titik tetap pada \textit{pipeline} komputasi terbuka R menghasilkan estimasi korelasi populasi terkoreksi $\bar{R} \approx 0,425$ ($S_Z \approx 67,72$; interval \textit{bootstrap} 95\% untuk $\bar{R}$ $[0,407; 0,442]$). Estimasi ini peka terhadap simpangan baku nasional tidak resmi, karena perubahan seluruh simpangan baku sebesar $\pm 10\%$ menggeser $\bar{R}$ menjadi $0,300$--$0,519$. Distribusi skor nasional juga menyimpang dari normal (miring ke kanan dengan ekor kiri terpotong), dan simulasi menunjukkan bahwa asumsi normal mengubah persentil terutama di ekor distribusi. \textit{Confirmatory factor analysis} (CFA) menunjukkan bahwa model dua faktor (Literasi dan Numerasi) lebih sesuai daripada model dengan satu \textit{g-factor} (\textit{Comparative Fit Index} $0,967$ berbanding $0,932$; $\Delta\chi^2 = 91,55$, $p < 0,001$; korelasi antarfaktor $\phi = 0,813$). Hasil analisis juga mengindikasikan bahwa bias audiens platform tidak dapat diatasi semata-mata melalui \textit{framing} pengumpulan data yang lebih inklusif.

\textit{Kata kunci:} UTBK-SNBT, korelasi antarsubtes, psikometri pendidikan, bias seleksi, data \textit{crowdsourced}, Thorndike Case III, analisis sensitivitas

\newpage

\begin{otherlanguage}{english}
\section*{Abstract}
\noindent
The inter-subtest correlation matrix of the 2026 Indonesian university entrance examination (UTBK-SNBT) is not publicly available, and the subtest-level descriptive statistics in circulation derive from an unofficial, unverified internal table attributed to the national admissions board (SNPMB). This study aims to estimate the correlation structure across the eight subtests of the 2026 UTBK-SNBT using crowdsourced score data from platform X ($N = 1,886$), whilst identifying the characteristics of methodological bias inherent in this approach and the implications for national ranking estimation. The observed raw off-diagonal mean correlation was $0.300$, which is moderate-to-low for a set of interrelated cognitive ability tests. Comparison with these unofficial national parameters indicates a selection bias towards higher scores, consistent with restriction of range. A Thorndike Case III adjustment via fixed-point iteration within an open-source R computational pipeline yielded a corrected population correlation of $\bar{R} \approx 0.425$ ($S_Z \approx 67.72$; 95\% bootstrap interval for $\bar{R}$ $[0.407, 0.442]$). This estimate is sensitive to the unofficial national standard deviations, as changing all of them by $\pm 10\%$ shifts $\bar{R}$ to between $0.300$ and $0.519$. The national score distribution also departs from normality (right-skewed with a truncated left tail), and simulation shows that the normality assumption mainly alters percentiles in the distribution tails. A confirmatory factor analysis showed that a two-factor model (Literacy and Numeracy) fit better than a single general-factor model (CFI $0.967$ versus $0.932$; $\Delta\chi^2 = 91.55$, $p < 0.001$; $\phi = 0.813$). Furthermore, the results indicate that platform-audience bias cannot be resolved solely through more inclusive data-collection framing.

\textit{Keywords:} UTBK-SNBT, inter-subtest correlation, educational psychometrics, selection bias, crowdsourced data, Thorndike Case III, sensitivity analysis
\end{otherlanguage}